\section*{[2024 EXAM] Problem 1: Weak form, Neumann BCs, Lax--Milgram}
We are given $u \in H^2(\Omega)$ solving
\[
  \begin{cases}
    - \Delta u = f       & \text{in } \Omega \subset \mathbb{R}^3, \\
    \nabla u \cdot n = 0 & \text{on } \partial\Omega,
  \end{cases}
\]
for $f \in L^2(\Omega)$, where $n$ is the unit outward normal.
\subsection*{(a) Weak formulation, bilinear form $a$ and functional $g$}
We want to write the problem in the abstract form
\[
  a(u,v) = g(v) \qquad \forall v \in V,
\]
by identifying a bilinear form $a(\cdot,\cdot)$ and a bounded linear functional $g$.

\textbf{Choice of space.}
The natural energy space for Neumann problems is $H^1(\Omega)$.
(Strictly speaking, the solution is only determined up to a constant, so one should work
in $V = H^1(\Omega)/\mathbb{R}$ or in the subspace of mean-zero functions.
Following the exam note, we will write $V = H^1(\Omega)$ and ignore this kernel issue in (a).)

\textbf{Derivation of the weak form.}
Let $v \in C^\infty(\overline{\Omega})$ (later we take $v \in H^1(\Omega)$ by density).
Multiply the PDE by $v$ and integrate over $\Omega$:
\[
  \int_\Omega (-\Delta u)\, v \,dx = \int_\Omega f\,v\,dx.
\]
Integrate by parts in $\Omega$:
\[
  \int_\Omega (-\Delta u)\, v \,dx
  = \int_\Omega \nabla u \cdot \nabla v \,dx
  - \int_{\partial\Omega} (\nabla u \cdot n)\, v \,ds.
\]
Using the Neumann boundary condition $\nabla u \cdot n = 0$ on $\partial\Omega$, we obtain
\[
  \int_\Omega (-\Delta u)\, v \,dx
  = \int_\Omega \nabla u \cdot \nabla v \,dx.
\]
Hence the weak form is:
\[
  \int_\Omega \nabla u \cdot \nabla v \,dx
  = \int_\Omega f\,v\,dx
  \quad\forall v \in H^1(\Omega).
\]
We now define
\[
  a(u,v) := \int_\Omega \nabla u \cdot \nabla v \,dx,\qquad
  g(v) := \int_\Omega f\,v\,dx,
\]
and we take $V = H^1(\Omega)$ (or $H^1(\Omega)/\mathbb{R}$ for full rigor).

\textbf{Boundedness of $g$.}
For $v \in H^1(\Omega)$, by Cauchy--Schwarz,
\[
  |g(v)| = \left|\int_\Omega f\,v\,dx\right|
  \le \|f\|_{L^2(\Omega)}\,\|v\|_{L^2(\Omega)}
  \le \|f\|_{L^2(\Omega)}\,\|v\|_{H^1(\Omega)},
\]
so $g$ is a bounded linear functional on $V$.
Hence the Poisson problem can be written as:
\[
  \boxed{
    \text{Find } u \in V \text{ s.t. } a(u,v) = g(v)\quad\forall v \in V,
  }
\]
with
\[
  a(u,v) = \int_\Omega \nabla u \cdot \nabla v\,dx,\qquad
  g(v) = \int_\Omega f\,v\,dx,\qquad V = H^1(\Omega)
  \text{ (up to constants)}.
\]

\subsection*{(b) V-ellipticity \& uniqueness (Lax--Milgram)}

\textbf{Definition (V-elliptic).}
Let $V$ be a Hilbert space with norm $\|\cdot\|_V$ and $a : V\times V \to \mathbb{R}$ a bilinear form.
We say that $a$ is \emph{$V$-elliptic} (or \emph{coercive on $V$}) if there exists a constant $\alpha>0$ such that
\[
  a(v,v) \;\ge\; \alpha\,\|v\|_V^2 \quad \forall v \in V.
\]
\textbf{Boundedness of $a$.}
For $u,v \in V$,
\[
  |a(u,v)|
  = \left|\int_\Omega \nabla u \cdot \nabla v\,dx\right|
  \le \|\nabla u\|_{L^2(\Omega)}\,\|\nabla v\|_{L^2(\Omega)}
  \le \|u\|_{H^1(\Omega)}\,\|v\|_{H^1(\Omega)},
\]
so $a$ is bounded on $V\times V$.
\textbf{Coercivity of $a$.}
We equip $V = H^1(\Omega)$ with the standard norm
\[
  \|v\|_{H^1(\Omega)}^2 := \|v\|_{L^2(\Omega)}^2 + \|\nabla v\|_{L^2(\Omega)}^2.
\]
We are given the Poincaré inequality (formally)
\[
  \|v\|_{L^2(\Omega)} \le C_\Omega \|\nabla v\|_{L^2(\Omega)} \quad\forall v\in V,
\]
so
\[
  \|v\|_{H^1(\Omega)}^2
  = \|v\|_{L^2(\Omega)}^2 + \|\nabla v\|_{L^2(\Omega)}^2
  \le C_\Omega^2 \|\nabla v\|_{L^2(\Omega)}^2 + \|\nabla v\|_{L^2(\Omega)}^2
  = (1+C_\Omega^2)\,\|\nabla v\|_{L^2(\Omega)}^2.
\]
But
\[
  a(v,v) = \int_\Omega |\nabla v|^2 \,dx = \|\nabla v\|_{L^2(\Omega)}^2,
\]
so we can rewrite the previous inequality as
\[
  \|v\|_{H^1(\Omega)}^2 \le (1+C_\Omega^2)\,a(v,v).
\]
Equivalently,
\[
  a(v,v) \;\ge\; \frac{1}{1+C_\Omega^2}\,\|v\|_{H^1(\Omega)}^2
  =:\alpha\,\|v\|_V^2,\quad \alpha = \frac{1}{1+C_\Omega^2}>0.
\]
Thus $a$ is $V$-elliptic.

\emph{Remark on the kernel:}
For a pure Neumann problem, constant functions satisfy $a(c,c) = 0$ but $c\neq 0$, so $a$ is not coercive on all of $H^1(\Omega)$.
The proper setting is $V = H^1(\Omega)/\mathbb{R}$ or
\[
  V = \bigl\{ v\in H^1(\Omega) : \int_\Omega v\,dx = 0\bigr\},
\]
where Poincaré inequality holds without the constant functions, and the above coercivity argument is rigorous.
The exam note allows us to ignore this subtlety and write $V=H^1(\Omega)$.
\textbf{Boundedness of $g$ revisited.}
From (a) we already have
\[
  |g(v)| \le \|f\|_{L^2(\Omega)}\|v\|_{L^2(\Omega)}
  \le \|f\|_{L^2(\Omega)}\,C_\Omega\,\|\nabla v\|_{L^2(\Omega)}
  \le \|f\|_{L^2(\Omega)}\,C_\Omega\,\|v\|_{H^1(\Omega)},
\]
so $g\in V^\ast$ is bounded.
\textbf{Application of Lax--Milgram.}
We have shown:
\begin{itemize}
  \item $a(\cdot,\cdot)$ is bilinear and bounded on $V\times V$.
  \item $a$ is $V$-elliptic: $\exists \alpha>0$ such that $a(v,v)\ge\alpha\|v\|_V^2$ for all $v\in V$.
  \item $g\in V^\ast$ is a bounded linear functional.
\end{itemize}
Therefore, by the Lax--Milgram theorem, there exists a unique $u\in V$ such that
\[
  a(u,v) = g(v) \qquad \forall v\in V.
\]
Concretely, this says:
There exists a unique
\[
  \boxed{
    u \in V \text{ solving }
    \int_\Omega \nabla u \cdot \nabla v\,dx = \int_\Omega f\,v\,dx
    \quad\forall v\in V,
  }
\]
where, for full rigor, $V = H^1(\Omega)/\mathbb{R}$ (or equivalently, the subspace of $H^1(\Omega)$ of mean-zero functions).
Thus the weak form derived in (a) is well posed.
